% !PRINTED 348
\weberpage{348}
%<<CONT
\[
1, \quad 4a_1^2 - 10a_0a_2, \quad 6a_2^2 - 12a_1a_3
\]
\[
6a_3^2 - 12a_2a_4, \quad 4a_4^2 - 10a_3a_5, \quad 1
\]
bilden, so erhalten wir die Vorzeichen
\[
+ \quad + \quad - \quad + \quad - \quad +
\]
und (8) hat also nach \S.\,101, Lehrsatz VIII vier imaginäre Wurzeln. Die positive Wurzel ist jedenfalls grösser als 1, da die linke Seite von (8) für $x = 0$ und $x = 1$ negativ ist.

Um nun zu entscheiden, welche der Wurzeln den grössten absoluten Werth hat, wollen wir in (8) $x = 1 : z$ setzen und dann nach \S.\,96, 97 untersuchen, wie viele Wurzeln im Inneren des Einheitskreises in der $z$-Ebene liegen. Setzen wir also in
\[
z^5 + 8z^4 + 36z^3 - z^2 + 82z - 1
\]
$z = \cos\vartheta + i\sin\vartheta$, so erhalten wir den reellen und imaginären Theil:
\[
\varphi = \cos 5\vartheta + 8\cos 4\vartheta + 36\cos 3\vartheta - \cos 2\vartheta + 82\cos\vartheta - 1
\]
\[
\psi = \sin 5\vartheta + 8\sin 4\vartheta + 36\sin 3\vartheta - \sin 2\vartheta + 82\sin\vartheta.
\]
Es lässt sich nun zeigen, dass $\psi$ positiv ist, wenn $0 < \vartheta < \pi$ und folglich negativ, wenn $0 > \vartheta > -\pi$, dass also der Einheitskreis in zwei Segmente getheilt ist, in denen $\psi$ entgegengesetzte Vorzeichen hat. Daraus folgt nach \S.\,96, dass im Inneren des Kreises nur eine Wurzel liegen kann, und dies ist die reelle Wurzel, während die imaginären Wurzeln ausserhalb liegen. Es folgt dann daraus für die Gleichung (8), dass die reelle Wurzel den grössten absoluten Werth hat, und dass also dieser durch die \gesperrt{Gräffe'sche Methode} gefunden wird.

Um nun diese Eigenschaft der Function $\psi$ nachzuweisen, setzen wir $2\cos\vartheta = \xi$, und wenden die goniometrischen Formeln an:
\[
\begin{aligned}
\sin 5\vartheta &= \sin\vartheta(\xi^4 - 3\xi^2 + 1) \\
\sin 4\vartheta &= \sin\vartheta(\xi^3 - 2\xi) \\
\sin 3\vartheta &= \sin\vartheta(\xi^2 - 1) \\
\sin 2\vartheta &= \sin\vartheta\,\xi.
\end{aligned}
\]
Dadurch erhalten wir
\[
\psi = \sin\vartheta[\xi^4 + \xi^2(8\xi + 33) + (47 - 17\xi)],
\]
woraus man sofort sieht, dass der Factor von $\sin\vartheta$ positiv bleibt, so lange $\xi$ zwischen $-2$ und $+2$ liegt.
